% HD-I20. GPT-6 Astra Pro, /root/imprimitive, inherited effort.
% 21 November 2026. Exact cotangent trace, shared-value map and new strata.
\subsection{HD-I20: the exact quadratic trace and ramified uncolored cycles}
\label{sec:HD-I20}

Use the actual quotient tower $X\to Y\xrightarrow{h}\mathbf P^1$
of HD-I01, with its faithful primitive character and balanced
signature $(n,n)$. Throughout, $d=3,4,6$, $e\mid d$, $n\geq5$,
and $m\geq\lceil(n^2-2n+3)/2\rceil$. Write $g=g(Y)$ and
let $R$ be the reduced colored divisor on $Y$. Let $b$ be the
number of branch values, $\nu$ the number of nonreflection values,
$B$ their reduced divisor, and $A=b-\nu$ the number of reflection
values. Let $Q$ contain once every uncolored permutation cycle
over $B$, including unramified ones. The line bundle of HD-I19 is
\begin{equation}
 F=h^*T_{\mathbf P^1}(-h^*B+Q),\qquad \deg F=A-2n.
                                                        \tag{I20.1}
\end{equation}
Here and below $h^*T_{\mathbf P^1}(-h^*B+Q)$ means
$h^*(T_{\mathbf P^1}(-B))\otimes\mathcal O_Y(Q)$.

The formulas below compute the differential on every individual
Nielsen component at every point. They retain the actual trace
kernel; no statement that this kernel has a fixed dimension on
all branch strata is needed or asserted.

\paragraph{1. The shared-value correction is a specific cohomology map.}
Let $D_{\mathrm{jet}}$ be the divisor of HD-I16 and put
$E=h^*T_{\mathbf P^1}$. Thus
$D_{\mathrm{jet}}=h^*B-Q$ and $F=E(-D_{\mathrm{jet}})$.
For arbitrary $\nu$, the local constraint quotient
$H^0(E|_{D_{\mathrm{jet}}})\to\mathbf C^N$ has kernel
\[
 U=\bigoplus_{i=1}^{\nu}T_{b_i}\mathbf P^1,
       \qquad N=2m+2n-b.
\]
The summand at $b_i$ specifies the common constant coefficient
on every relevant local jet over that value. Its other coefficients
are zero. Since $H^1(E)=0$ in the stated degree range, the
restriction sequence gives a connecting map
\[
 \delta:H^0(E|_{D_{\mathrm{jet}}})\longrightarrow H^1(Y,F).
\]
Its restriction to $U$ factors as the following actual maps:
\begin{equation}
 U\xrightarrow{\delta_0}
 H^1\bigl(\mathbf P^1,T_{\mathbf P^1}(-B)\bigr)
 \xrightarrow{\ \kappa\ }H^1(Y,F).                    \tag{I20.2}
\end{equation}
Here $\delta_0$ is the connecting map for target evaluation;
$\kappa$ is pullback in cohomology, followed by the inclusion
$h^*T_{\mathbf P^1}(-h^*B)\hookrightarrow F$.

For completeness, pull back the target evaluation sequence, and
map it to
$0\to F\to E\to E|_{D_{\mathrm{jet}}}\to0$ by the displayed
inclusion, the identity on $E$, and restriction from the full
fibres $h^*B$ to $D_{\mathrm{jet}}$. The target constant at
$b_i$ pulls back to the same constant on each local jet.
Naturality of the connecting maps proves (I20.2), including
the marked colored points and the shortened uncolored jets.
Also $\delta_0$ is onto, since $H^1(T_{\mathbf P^1})=0$.

Write $r_\kappa=\operatorname{rank}\kappa$. The image of
$\ker J$ in $U$ is exactly $\ker(\delta|_U)$, so
\[
 \dim\ker J=h^0(F)+\nu-r_\kappa.
\]
Quotienting by the three pulled-back target vector fields as
in HD-I16 therefore proves
\begin{equation}
 \boxed{\operatorname{rank}d\mathcal F
       =b-\nu-h^0(F)+r_\kappa
       =2n+g-1-h^1(F)+r_\kappa.}                     \tag{I20.3}
\end{equation}
Thus the additional shared-value map has a domain of dimension
$\max(0,\nu-3)$, regardless of the much larger number $N$ of
original local constraints. For $\nu\leq3$ it vanishes and
(I20.3) is exactly HD-I19.

\paragraph{2. Its dual is the restricted quadratic trace.}
Let $R_h$ be the ramification divisor of $h$. Define
$D_{\mathrm{ref}}$ as the sum of the unique relevant points
over the reflection values. Such a point is either an uncolored
ramification point of index two or a colored unramified point.
Every summand has multiplicity one. The local jet lengths give
the exact divisor identity
\begin{equation}
 D_{\mathrm{jet}}=R_h+R-D_{\mathrm{ref}},
 \qquad F\simeq T_Y(-R+D_{\mathrm{ref}}).              \tag{I20.4}
\end{equation}
The second isomorphism is induced by the differential of the
actual map $h$: $dh:T_Y\to E$ has divisor $R_h$.
Consequently Serre duality identifies
\[
 H^1(F)^*=H^0\bigl(Y,\omega_Y^2(R-D_{\mathrm{ref}})\bigr).
\]
Dualizing the specified pullback and inclusion in (I20.2)
gives the canonical-sheaf trace, with projection formula.
Under (I20.4) this is exactly
\begin{equation}
 T_0:H^0\bigl(Y,\omega_Y^2(R-D_{\mathrm{ref}})\bigr)
       \longrightarrow
       H^0\bigl(\mathbf P^1,\omega_{\mathbf P^1}^2(B)\bigr),
 \qquad
 \Omega\longmapsto
       \sum_{h(z)=w}\frac{\Omega(z)}{(h'(z))^2}.
                                                        \tag{I20.5}
\end{equation}
The displayed expression includes the factor $dw^2$ after
writing $\Omega(z)=f(z)\,dz^2$. It is independent of coordinates:
divide by one copy of $dh$, take the ordinary trace of a
canonical differential, and retain the remaining target
canonical factor. This describes the precise morphism dual
to $\kappa$, rather than only a map of the same dimensions.
In particular $\operatorname{rank}T_0=r_\kappa$ and
\begin{equation}
 \boxed{\operatorname{rank}d\mathcal F
             =2n+g-1-\dim\ker T_0.}                  \tag{I20.6}
\end{equation}

Here is a local verification of the poles and coefficients.
For $w=z^\ell$ and
$\Omega=\sum_k a_k z^k\,dz^2$, summing over the $\ell$
inverse branches cancels terms unless
$\ell\mid k-2\ell+2$. The term with a simple pole downstairs
has $k=\ell-2$, and its coefficient is
\begin{equation}
                    \frac{a_{\ell-2}}{\ell}.
                                                        \tag{I20.7}
\end{equation}
At a colored point $k\geq-1$; at an uncolored point $k\geq0$.
These bounds allow no pole of order two or more after the
sum. An uncolored unramified point gives no pole, a colored
unramified point contributes $a_{-1}$, and an uncolored
simple ramification point contributes $a_0/2$.
At a reflection value the last two cases exhaust the relevant
points, and (I20.4) respectively removes the pole or requires
a zero. Thus (I20.5) has poles only at $B$ as claimed.

\paragraph{3. The full cotangent and period maps.}
Let $B_{\mathrm{all}}$ be the reduced divisor of all $b$ branch
values. The same formula defines
\begin{equation}
 T_h:H^0\bigl(Y,\omega_Y^2(R)\bigr)
       \longrightarrow
       H^0\bigl(\mathbf P^1,
                         \omega_{\mathbf P^1}^2(B_{\mathrm{all}})\bigr).
                                                        \tag{I20.8}
\end{equation}
This is the actual cotangent differential of the forgetful
map $\mathcal F$. To check the deformation identification,
move a branch value by $\epsilon\beta$ and use the local
coordinate $w=z^\ell$. The change from a local trivialization
around its preimage to the trivialization away from that
preimage is represented by
\[
            \frac{\beta}{\ell z^{\ell-1}}\,\partial_z.
\]
Take the Cech differential as this local trivialization minus
the outside one. Its Serre pairing with $\Omega$ is the
residue of $\beta\Omega/(\ell z^{\ell-1}dz)$, namely
$\beta a_{\ell-2}/\ell$. At a colored unramified point this
is the motion of the marked point in the deformation space
of $(Y,R)$. At an uncolored ramification point the centre is
unmarked, exactly as in HD-I16. Summing at all points over a
common value gives (I20.7). These residues are also the
standard cotangent pairing with the motion of the target
branch value. They prove (I20.8) with the indicated Cech
convention, and automatically annihilate target vector fields.

If $T_h\Omega=0$, its pole coefficient at each reflection
value is zero. The unique relevant point then forces the
same zero or pole removal used in (I20.4). Hence
\[
 \ker T_h\subset
        H^0\bigl(\omega_Y^2(R-D_{\mathrm{ref}})\bigr),
 \qquad \ker T_h=\ker T_0.                            \tag{I20.9}
\]
Conversely a section in the latter kernel has zero full trace,
so the equality is an equality of actual subspaces.

Let $\mu$ be the primitive eigendifferential multiplication
of HD-I05. The period codifferential for the original source,
under the fixed polarized isogeny of HD-I02, is the composite
\begin{equation}
 H^0(\omega_YL_\chi)\otimes H^0(\omega_YL_{\bar\chi})
       \xrightarrow{\ \mu\ }H^0(\omega_Y^2(R))
       \xrightarrow{\ T_h\ }
       H^0\bigl(\omega_{\mathbf P^1}^2(B_{\mathrm{all}})\bigr).
                                                        \tag{I20.10}
\end{equation}
This follows by composing the actual cyclic period morphism
with $\mathcal F$; the finite polarized isogeny identifies
their rational variations and changes coordinates invertibly.
Thus every individual period rank is exactly
$\operatorname{rank}(T_h\mu)$. For $g\leq n-2$, HD-I12
proves $\mu$ onto at every point, and (I20.6) itself is the
exact period rank. In the three larger genera (I20.10)
specifies the additional intersection with the cyclic
multiplication image without assuming its restriction is
generic along the given Hurwitz component.

\paragraph{4. A new class with a completely evaluated shared-value map.}
Suppose every uncolored cycle over a nonreflection value has
length at least two. This is an explicit branch-profile
condition; it permits arbitrarily many such values and colored
cycles of any length. The ordinary function trace extends to
\[
           \operatorname{Tr}_h:h_*\mathcal O_Y(Q)
                                      \longrightarrow\mathcal O_{\mathbf P^1}.
\]
Indeed the only new allowed local term is a simple pole
$z^{-1}$ at a point of $Q$. In the coordinate $w=z^\ell$,
its trace is zero for $\ell\geq2$. All nonnegative powers
have regular trace. Away from these points there is no pole.
The composite of the inclusion of constants and this trace
is multiplication by $m$. Therefore $(1/m)\operatorname{Tr}_h$
is a retraction of
$\mathcal O_{\mathbf P^1}\hookrightarrow h_*\mathcal O_Y(Q)$.

Tensoring this actual splitting with $T_{\mathbf P^1}(-B)$
and using the projection formula shows that $\kappa$ is
injective. For $\nu\geq3$ its rank is $\nu-3$, and hence
\begin{equation}
 \boxed{\operatorname{rank}d\mathcal F=b-3-h^0(Y,F)
 \quad\text{if }\nu\geq3\text{ and every uncolored cycle over }B
                       \text{ is ramified}.}          \tag{I20.11}
\end{equation}
In particular, $A<2n$ makes $\deg F<0$, so the forgetful
map is immersive of rank $b-3$ at every point of every such
component. At $A=2n$, its rank is $b-4$ when $F$ is trivial
and $b-3$ otherwise. All these are period ranks too when
$g\leq n-2$. Thus (I20.11) evaluates the shared-value map
for a further entire branch-profile class, beyond $\nu\leq3$.

\paragraph{5. Actual four-nonreflection components.}
The preceding class is nonempty for every proper divisor
$e<d$, every $n\geq5$, and every even $m$ in the prompt's
range. There are components whose period rank is everywhere
$2n-5$. For $e=d$, every $n\geq6$ and every such even $m$
admit components of everywhere rank $2n-7$.

Here are explicit generating tuples and their exact group
calculation. Write $m=2a$, and index the letters by
$o_j=2j+1$, $v_j=2j+2$ for $j\in\mathbf Z/a\mathbf Z$.
Let $C$ have the matching cycles $(o_j\,v_j)$; color its
first two cycles $1,-1$ and all others zero. Let $D$ have
the matching cycles $(v_j\,o_{j+1})$. When $e<d$, color its
first cycle $e$ and all others zero, and put $k=3$.
When $e=d$, color its first two cycles $1,-1$ and all others
zero, and put $k=4$. A colored two-cycle can be represented
by one nontrivial matrix entry and one entry equal to one;
its square has the specified color on both diagonal entries.
These choices define actual monomial matrices in $G(d,e,m)$.
Their total colors are respectively $0,e$ in the proper-divisor
case and $0,0$ in the case $e=d$.

Put $s=(o_0\,v_0)$ and $t=(v_0\,o_1)$, with zero colors,
and take
\begin{equation}
 \bigl(C,C^{-1},D,D^{-1},s,s,t,t,
          \underbrace{s,s,\ldots,s,s}_{n-k-2\ \mathrm{pairs}}\bigr).
                                                        \tag{I20.12}
\end{equation}
The exponent is nonnegative in the stated ranges. Its product
is one. The four matching matrices are nonreflections, and
all remaining entries are reflections.

To prove generation, the underlying permutation of $CD$ sends
$o_j$ to $o_{j-1}$ and $v_j$ to $v_{j+1}$. The conjugates of
$s$ and $t$ by its powers are colored involutory transpositions
whose underlying edges are respectively
\[
            (o_{-j},v_j),\qquad (v_j,o_{1-j}).
\]
Their union is connected: its two alternating edges connect
$o_i$ to $v_{-i}$ and then to $o_{i+1}$. Choose a spanning
tree. Assign a phase to its root, and then assign the phase
at each new vertex so that conjugation by the resulting
diagonal matrix removes the color on its tree edge. This
constructs an actual diagonal conjugacy under which the tree
reflections become the ordinary tree transpositions. They
generate $S_m$. In particular the original generated group
contains a diagonal conjugate of $S_m$, and its conjugation
action on diagonal matrices is the usual permutation action.

The diagonal matrix $C^2$ has values $1,-1,0$ on its first,
second and remaining matching pairs. The stated degree range
has $a\geq5$, so the value zero is present. Taking its
commutator with an element of the preceding symmetric group
that swaps a coordinate of value one and one of value zero
gives the unit diagonal difference on those two coordinates.
Conjugation then gives all unit differences. These generate
the full augmentation diagonal group
\[
        \{(x_1,\ldots,x_m)\in(\mathbf Z/d)^m:
                                          \textstyle\sum_i x_i=0\}.
\]
Together with the conjugate symmetric group this is exactly
the kernel of total color in $G(d,1,m)$: cancel an arbitrary
permutation first, then its remaining total-zero diagonal.
Finally $D$ has total color $e$ for $e<d$; for $e=d$ all
generators have total color zero. Taking the exact preimage
of $e\mathbf Z/d$ therefore gives $G(d,e,m)$ in every case.
This proves full monodromy, including its diagonal subgroup.

There are $2(n-k)$ reflection values and four nonreflection
values. The four matching permutations contribute ramification
$2m$ to $h$, and the reflections contribute $2(n-k)$.
Riemann--Hurwitz consequently gives
\[
       g=n-k+1,\qquad \deg R=2k,\qquad
       b=4+2(n-k),\qquad \deg F=-2k.
\]
For each colored cycle of $C$ or $D$, its inverse has the
complementary color. Each such pair contributes one to the
sum of fractional colors, so the two primitive eigenspace
dimensions are both $g-1+k=n$. Every uncolored matching
cycle has length two. Thus (I20.11) applies, $h^0(F)=0$,
and $g\leq n-2$ makes the cyclic multiplication onto.
The exact period rank is $b-3=2n-2k+1$, as asserted.
The product-one full-group tuples give actual connected
covers by the Riemann existence correspondence. No generic
group-generation or differential-rank assertion is being
inferred from the finite checks below.

\paragraph{Scope and checks.}
The author, GPT-6 Astra Pro with inherited effort, read this
complete argument on 21 November 2026. The morphisms checked
are the natural diagram giving $\kappa$, the differential
$dh$ identifying $F$, canonical trace and Serre duality, the
local Cech residue computation, and the actual function-trace
splitting. The four-value construction has a separate exact
monomial certificate. Formula (I20.10) is the general
componentwise rank gate; (I20.11) and (I20.12) evaluate it in
the stated additional strata. They do not assign a uniform
rank to all remaining arbitrary Nielsen components.
